% Generated by tools/build_v6.py; edit the assembly rules, not this file.
\section{Exact numerical certificates on finite laws}
\label{sec:computation}

The computations in this section produce two-sided certificates for
eight finite rational instances. Numerical conic optimization selects
candidate policies and dual priors. Separate rational-arithmetic
verifiers then check the resulting bounds, including the observation
budget and the stopping criterion. An approximate solver objective is
never used as a certified endpoint.

\subsection{What the two endpoints mean}

Write $U_N$ for the risk bound of a verified observable cap-$N$ policy,
and $L_N$ for the value of a verified feasible genie dual. Their roles are
\[
 L_N\le G_N\le P_\infty\le P_N\le U_N,
 \qquad L_N/8\le R_{\mu,\rho}(B)\le U_N.
\]
Thus $L_N$ is a lower bound on a program value. It is not, by itself,
the reported lower bound on fully adaptive minimax risk. We also compute
the stronger pair-specific adaptive-risk lower bound
$\underline r_N$ from
\eqref{eq:direct-genie-risk}--\eqref{eq:direct-genie-risk-squared}.
Its target separation is that of the absorbed midpoint pair.

The cap search stops only after the exact rational comparison
\begin{equation}\label{eq:computed-stopping-test}
 U_N^2\le4L_N^2
\end{equation}
has passed. This implies $U_N\le16R_{\mu,\rho}(B)$.
The comparison uses the stored rational squared values, not the rounded
decimals in the tables.

\subsection{The benchmark laws and results}

Every instance has $\rho=1/10$. The four laws are
\begin{align*}
 \mu_{\rm one}&=\delta_{1/4},\\
 \mu_{\rm three}&=\tfrac34\delta_{2/5}
                  +\tfrac15\delta_{1/5}+\tfrac1{20}\delta_{2/25},\\
 \mu_{\rm hard}&=\tfrac{99}{100}\delta_{9/20}
                  +\tfrac1{100}\delta_{2/25},\\
 \mu_{\rm label}&=\delta_{1/2}.
\end{align*}
Each was run at $B=100$ and $B=1000$. The candidate caps were
$4,8,16,32$, with the perfect-label cases starting at $N=1$.
Seven instances passed \eqref{eq:computed-stopping-test} at their first
cap. The three-scale law at $B=1000$ failed at $N=4$ and passed at $N=8$.
There were therefore nine attempted brackets and eight successful ones.
The failed bracket is retained in the source bundle; it remains a valid
bracket but does not meet the requested factor-two program criterion.

\begin{table}[htbp]
\centering
\small
\setlength{\tabcolsep}{4pt}
\begin{tabular}{lrr rrr}
\toprule
Law & $B$ & $N$ & $L_N$ & $\underline r_N$ & $U_N$\\
\midrule
One gap & 100 & 4 & 0.03474453 & 0.00619890 & 0.03891866\\
One gap & 1000 & 4 & 0.01461177 & 0.00196026 & 0.01825093\\
Three scales & 100 & 4 & 0.03331182 & 0.00573109 & 0.04551457\\
Three scales & 1000 & 8 & 0.01569090 & 0.00212025 & 0.02553208\\
Rare hard atom & 100 & 4 & 0.02912416 & 0.00459699 & 0.03009854\\
Rare hard atom & 1000 & 4 & 0.01105026 & 0.00145369 & 0.01166586\\
Perfect labels & 100 & 1 & 0.02572478 & 0.00384900 & 0.02572479\\
Perfect labels & 1000 & 1 & 0.00932054 & 0.00121716 & 0.00932055\\
\bottomrule
\end{tabular}
\caption{Successful exact brackets. Both lower columns are rounded
downward and the upper column upward. The adaptive minimax risk lies
in $[\underline r_N,U_N]$; the unrestricted flow-program value lies
in $[L_N,U_N]$. These two intervals have different meanings.}
\label{tab:exact-brackets}
\end{table}

For the failed three-scale attempt at $B=1000,N=4$, the verified
program endpoints, rounded outward, are $0.01359364$ and $0.02989647$,
and the exact ratio exceeds two. At $N=8$ the program ratio is about
$1.6272$, so the exact test passes. Across the eight successful
instances, the program bracket ratios range from essentially one to
$1.6272$. The sharper adaptive-risk bracket ratios range from about
$6.27$ to $12.05$, all below sixteen.

These are small benchmark instances. Their modest successful caps
provide no uniform bound on the cap or runtime needed for another law,
rarity, or budget. The convergence theorem is qualitative and the
certificates concern the particular rational inputs displayed above.

\subsection{An analytic check with perfect labels}

For $\mu=\delta_{1/2}$, the first observation reveals the sign exactly.
Further observations add no information. For every cap $N\ge1$,
\begin{equation}\label{eq:perfect-label-program}
 P_N(B)^2=P_\infty(B)^2=G_N(B)^2
 =\frac{ac}{a+c},
 \qquad
 a=\frac{\rho^2}{4},\quad c=\frac{\rho(1-\rho)}{B}.
\end{equation}
This is an analytic identity for the flow program, without an assertion
that it is the exact fully adaptive minimax risk.

To verify \eqref{eq:perfect-label-program}, conditional averaging removes
randomness in the output given the revealed sign and cannot increase
its second moment. Let $\lambda\le B$ be the active Poisson intensity,
let $y_0,y_1$ be the two outputs, and put
$s=\lambda(y_1-y_0)$. Optimizing the deterministic offset gives worst
bias $\rho|1-s|/2$. At the prior $q=\rho$, centering the outputs minimizes
the raw second moment, giving the variance lower bound
$s^2\rho(1-\rho)/\lambda$. It is attained by
\[
 y_0=-\rho s/\lambda,\qquad y_1=(1-\rho)s/\lambda.
\]
Since $\rho\le1/4$, their second moment is increasing in $q$ on
$[0,\rho]$, so its worst value occurs at $\rho$. Choosing $\lambda=B$
therefore reduces the objective exactly to
$a(1-s)^2+cs^2$, whose minimum is $ac/(a+c)$.
The genie already reveals nothing new after the first observation,
which proves the same equality for $G_N$.

For the two stored perfect-label instances, the exact rational analytic
value lies between the exact rational lower and upper certificates.
This was checked directly, without rounded decimal comparisons.

\subsection{How the certificates are repaired and checked}

The conic solver may return a status indicating limited numerical
accuracy. Such a status neither establishes nor invalidates a certificate.
The following separate checks determine feasibility.

For a lower certificate, the selected priors are quantized to rational
numbers, clipped to $[0,1]$, and, if needed, scaled to satisfy the rarity
constraint exactly. Variance and work references are regularized by
mixing a weight $10^{-8}$ of the constant prior $\rho/2$.
All binomial kernels and latent rewards are then recomputed exactly.
A numerical bisection proposes the work price. A new backward Bellman
pass uses rational arithmetic and rounds each potential upward to a
dyadic grid of denominator $2^{96}$. The price is increased if the
exact root inequality does not hold. An independent verifier reconstructs
the kernels and checks every rarity, stopping, continuation, terminal,
and root inequality.

For an upper certificate, a separate verifier works with the actual
observable policy: its stopping probabilities at each count state,
terminal outputs, Poisson intensity, and deterministic offset. It
computes exact conditional means, second moments, and observation
counts for every atom and sign. Total stopping probability must equal
one exactly. The worst-prior linear expressions are evaluated by an
exact fractional-knapsack calculation over $\mathcal Q_\rho$.
If rationalization raises expected work above $B$, the intensity is
reduced by the exact required factor and outputs are increased by its
reciprocal. This preserves means, and the variance bound is recomputed.
The reported risk upper is
\[
 U_N^2=
 \left(\sup_{q\in\mathcal Q_\rho}
                  |\operatorname{Bias}_q|\right)^2+
 \sup_{q\in\mathcal Q_\rho}\operatorname{Var}_q.
\]
The two maxima need not occur at the same prior; their sum remains a
valid upper bound. Integer square-root calculations round every
displayed lower RMS downward and every displayed upper RMS upward.
The verifier also checks the direction of the stored decimal rounding.

In addition to the optimized genie brackets, the bundle includes eight
ordinary finite-cap dual checks and sixteen fixed-quartet midpoint-genie
checks. Standard-library-only re-verification accepted all 33 stored
lower certificates and nine upper policies. Six targeted corruptions
were rejected: a rarity violation, an insufficient Bellman price, an
overstated lower value, an insufficient genie terminal reward, excessive
actual policy work, and an understated upper risk. These checks remain
active when Python is invoked with its optimization flag.

\subsection{Standalone reproduction}

The repository directory \nolinkurl{paper/rare-event/v2/checks/} is self-contained. It includes
a clearly attributed copy of the unchanged round-three primal solver,
the independent dual solver, the exact verifiers, and all certificate
files. It does not import a solver from a neighboring research directory.
From that directory, verification needs only Python's standard library:
\begin{verbatim}
python self_check_certificates.py
python verify_finite_dual.py \
  optimized_genie/three_scales_N8_B1000_lower.json
python verify_policy_upper.py \
  optimized_genie/three_scales_N8_B1000_upper.json
\end{verbatim}
The numerical search used Python 3.12 with CVXPY 1.9.3, NumPy 2.5.4,
SciPy 1.18.1, and Clarabel 0.11.1. These observed versions are pinned
in \texttt{requirements.txt}. The principal search is reproduced by
\begin{verbatim}
python -m pip install -r requirements.txt
python solve_genie_bracket.py
\end{verbatim}
Re-running a numerical solver may select different feasible priors
and slightly different bounds. The mathematical guarantee comes from
the independently checked rational artifacts.
